%  A simple AAU report template.
\documentclass[11pt,oneside,a4paper,openright]{report}
\usepackage{wrapfig}
\usepackage{subcaption}
\usepackage[inline]{enumitem}
\usepackage[export]{adjustbox}
\usepackage{graphicx}
\usepackage{rotating}
\usepackage{multicol}
\usepackage{multirow}
\usepackage{pdfpages}
\usepackage{colortbl}
\usepackage{gensymb}
\usepackage{mathrsfs}
\usepackage{titlesec}
\usepackage{parskip}
\usepackage{enumitem}
\usepackage{booktabs}
\usepackage{xcolor}
\usepackage{mdframed}
\usepackage{hyperref}
\usepackage[margin=2.5cm]{geometry}
\usepackage{amsmath, amssymb}
\usepackage{booktabs}
\usepackage{parskip}
\usepackage{xcolor}
\usepackage{mdframed}
\usepackage{titlesec}
\usepackage{enumitem}
\usepackage[most]{tcolorbox}
% --- Summary box ---
\newmdenv[
  backgroundcolor=aaublue!8,
  linecolor=aaublue,
  linewidth=1.5pt,
  roundcorner=4pt,
  innerleftmargin=12pt,
  innerrightmargin=12pt,
  innertopmargin=10pt,
  innerbottommargin=10pt,
  skipabove=10pt,
  skipbelow=10pt
]{summarybox}

\newtcolorbox{keybox}{
    colback=green!5!white,
    colframe=green!50!black,
    fonttitle=\bfseries,
    title=Key Takeaway,
    rounded corners,
    left=6pt, right=6pt, top=4pt, bottom=4pt
}
 
% --- Note box ---
\newmdenv[
  backgroundcolor=lightgray,
  linecolor=accentgray,
  linewidth=1pt,
  roundcorner=4pt,
  innerleftmargin=12pt,
  innerrightmargin=12pt,
  innertopmargin=8pt,
  innerbottommargin=8pt,
  skipabove=8pt,
  skipbelow=8pt
]{notebox}

% --- Warning / key result box ---
\newmdenv[
  backgroundcolor=warnred!8,
  linecolor=warnred,
  linewidth=1.5pt,
  roundcorner=4pt,
  innerleftmargin=12pt,
  innerrightmargin=12pt,
  innertopmargin=10pt,
  innerbottommargin=10pt,
  skipabove=10pt,
  skipbelow=10pt
]{warnbox}
 
% --- Green result box ---
\newmdenv[
  backgroundcolor=greenbox,
  linecolor=greenborder,
  linewidth=1.2pt,
  roundcorner=4pt,
  innerleftmargin=12pt,
  innerrightmargin=12pt,
  innertopmargin=8pt,
  innerbottommargin=8pt,
  skipabove=8pt,
  skipbelow=8pt
]{resultbox}
% Insight box
\newmdenv[
  backgroundcolor=lightgray,
  linecolor=aaublue,
  linewidth=1.5pt,
  leftline=true, rightline=false, topline=false, bottomline=false,
  innerleftmargin=10pt,
  innerrightmargin=10pt,
  innertopmargin=6pt,
  innerbottommargin=6pt,
]{insightbox}

\usepackage{listings}

\definecolor{aaublue}{RGB}{33, 83, 130}
\definecolor{lightgray}{RGB}{245, 245, 245}

\titleformat{\section}{\large\bfseries\color{aaublue}}{}{0em}{}[\titlerule]
\titleformat{\subsection}{\normalsize\bfseries}{}{0em}{}

\hypersetup{colorlinks=true, linkcolor=aaublue, urlcolor=aaublue}

\title{\textbf{Networks and Systems -- MM1}\\
\large Topic 1: System Definitions and Examples\\[0.5em]
\normalsize Lecture Notes Summary}
\author{Based on lectures by Hans-Peter Schwefel}
\date{AAU, February 2025}


\newcommand{\mhat}{\hat{m}}
\newcommand{\eps}{\epsilon}
% --- Colour definitions ---
\definecolor{aaublue}{RGB}{33, 87, 152}
\definecolor{lightgray}{RGB}{245, 245, 245}
\definecolor{accentgray}{RGB}{100, 100, 100}
\definecolor{warnred}{RGB}{180, 40, 40}
\definecolor{greenbox}{RGB}{230, 245, 230}
\definecolor{greenborder}{RGB}{60, 140, 60}

%%%%%%%%%%%%%%%%%%%%%%%%%%%%%%%%%%%%%%%%%%%%%%%%
% Language, Encoding and Fonts
% http://en.wikibooks.org/wiki/LaTeX/Internationalization
%%%%%%%%%%%%%%%%%%%%%%%%%%%%%%%%%%%%%%%%%%%%%%%%
% Select encoding of your inputs. Depends on
% your operating system and its default input
% encoding. Typically, you should use
%   Linux  : utf8 (most modern Linux distributions)
%            latin1 
%   Windows: ansinew
%            latin1 (works in most cases)
%   Mac    : applemac
% Notice that you can manually change the input
% encoding of your files by selecting "save as"
% an select the desired input encoding. 
\usepackage[utf8]{inputenc}
% Make latex understand and use the typographic
% rules of the language used in the document.
\usepackage[english]{babel}
% Use the palatino font
\usepackage[sc]{mathpazo}
\linespread{1.05}         % Palatino needs more leading (space between lines)
% Choose the font encoding
\usepackage[T1]{fontenc}
%%%%%%%%%%%%%%%%%%%%%%%%%%%%%%%%%%%%%%%%%%%%%%%%
% Graphics and Tables
% http://en.wikibooks.org/wiki/LaTeX/Importing_Graphics
% http://en.wikibooks.org/wiki/LaTeX/Tables
% http://en.wikibooks.org/wiki/LaTeX/Colors
%%%%%%%%%%%%%%%%%%%%%%%%%%%%%%%%%%%%%%%%%%%%%%%%
% load a colour package
\usepackage[table,xcdraw]{xcolor}
\usepackage{xcolor}
\definecolor{aaublue}{RGB}{33,26,82}% dark blue
% The standard graphics inclusion package
\usepackage{graphicx}
% Set up how figure and table captions are displayed
\usepackage{caption}
\captionsetup{%
  font=footnotesize,% set font size to footnotesize
  labelfont=bf % bold label (e.g., Figure 3.2) font
}
% Make the standard latex tables look so much better
\usepackage{array,booktabs}
% Enable the use of frames around, e.g., theorems
% The framed package is used in the example environment
\usepackage{framed}
\usepackage{threeparttable}

%%%%%%%%%%%%%%%%%%%%%%%%%%%%%%%%%%%%%%%%%%%%%%%%
% Mathematics
% http://en.wikibooks.org/wiki/LaTeX/Mathematics
%%%%%%%%%%%%%%%%%%%%%%%%%%%%%%%%%%%%%%%%%%%%%%%%
% Defines new environments such as equation,
% align and split 
\usepackage{amsmath}
% Adds new math symbols
\usepackage{amssymb}
% Use theorems in your document
% The ntheorem package is also used for the example environment
% When using thmmarks, amsmath must be an option as well. Otherwise \eqref doesn't work anymore.
\usepackage[framed,amsmath,thmmarks]{ntheorem}
%\usepackage[intoc, danish]{nomencl}
%\makenomenclature
%%%%%%%%%%%%%%%%%%%%%%%%%%%%%%%%%%%%%%%%%%%%%%%%
% Nomenclature
%%%%%%%%%%%%%%%%%%%%%%%%%%%%%%%%%%%%%%%%%%%%%%%%
\usepackage{etoolbox}
\usepackage{siunitx}
\usepackage{ragged2e}
\usepackage[intoc, danish]{nomencl}
\makenomenclature



%%%%%%%%%%%%%%%%%%%%%%%%%%%%%%%%%%%%%%%%%%%%%%%%
% Page Layout
% http://en.wikibooks.org/wiki/LaTeX/Page_Layout
%%%%%%%%%%%%%%%%%%%%%%%%%%%%%%%%%%%%%%%%%%%%%%%%
% Change margins, papersize, etc of the document
\usepackage[a4paper,top=3cm,bottom=2cm,left=3cm,right=3cm,marginparwidth=1.75cm]{geometry}
% Modify how \chapter, \section, etc. look
% The titlesec package is very configureable
\usepackage{titlesec, blindtext, color}
\definecolor{gray75}{gray}{0.75}
\newcommand{\hsp}{\hspace{20pt}}
\titleformat{\chapter}[hang]{\Huge\bfseries}{\thechapter\hsp\textcolor{gray75}{|}\hsp}{0pt}{\Huge\bfseries}
\titleformat*{\section}{\normalfont\Large\bfseries}
\titleformat*{\subsection}{\normalfont\large\bfseries}
\titleformat*{\subsubsection}{\normalfont\normalsize\bfseries}
%\titleformat*{\paragraph}{\normalfont\normalsize\bfseries}
%\titleformat*{\subparagraph}{\normalfont\normalsize\bfseries}

% Clear empty pages between chapters
\makeatletter
\def\clearemptydoublepage{%
  \clearpage
  {\pagestyle{empty}\if@twoside\ifodd\c@page\else
  \hbox{}\newpage
  \if@twocolumn\hbox{}\newpage\fi\fi\fi}
}
\makeatother

% Change the headers and footers
\usepackage{fancyhdr}
\pagestyle{fancy}
\fancyhf{} %delete everything
\rhead{\groupName}
\lhead{\rightmark}
\rfoot{Page \thepage\ of \pageref{LastPage}}
\renewcommand{\headrulewidth}{1pt}
\renewcommand{\footrulewidth}{1pt}

\fancypagestyle{plain}{%
\fancyhf{}% clear all header and footer fields
\rfoot{Page \thepage\ of \pageref{LastPage}} % except the center
\renewcommand{\headrulewidth}{0pt}%
\renewcommand{\footrulewidth}{1pt}
}

% Do not stretch the content of a page. Instead,
% insert white space at the bottom of the page
\raggedbottom
% Enable arithmetics with length. Useful when
% typesetting the layout.
\usepackage{calc}

%%%%%%%%%%%%%%%%%%%%%%%%%%%%%%%%%%%%%%%%%%%%%%%%
% Bibliography
% http://en.wikibooks.org/wiki/LaTeX/Bibliography_Management
%%%%%%%%%%%%%%%%%%%%%%%%%%%%%%%%%%%%%%%%%%%%%%%%
\usepackage[style=numeric-comp,backend=bibtex,bibencoding=utf8,sorting=none]{biblatex}
\addbibresource{bib/references}
%\usepackage{microtype}
%\usepackage{breakcites}
\usepackage{xurl}


%%%%%%%%%%%%%%%%%%%%%%%%%%%%%%%%%%%%%%%%%%%%%%%%
% Misc
%%%%%%%%%%%%%%%%%%%%%%%%%%%%%%%%%%%%%%%%%%%%%%%%
% Add bibliography and index to the table of
% contents
\usepackage[nottoc]{tocbibind}
% Add the command \pageref{LastPage} which refers to the
% page number of the last page
\usepackage{lastpage}
% Add todo notes in the margin of the document
\usepackage[
%  disable, %turn off todonotes
  colorinlistoftodos, %enable a coloured square in the list of todos
  textwidth=\marginparwidth, %set the width of the todonotes
  textsize=scriptsize, %size of the text in the todonotes
  ]{todonotes}

%%%%%%%%%%%%%%%%%%%%%%%%%%%%%%%%%%%%%%%%%%%%%%%%
% Hyperlinks
% http://en.wikibooks.org/wiki/LaTeX/Hyperlinks
%%%%%%%%%%%%%%%%%%%%%%%%%%%%%%%%%%%%%%%%%%%%%%%%
% Enable hyperlinks and insert info into the pdf
% file. Hypperref should be loaded as one of the 
% last packages
\usepackage{hyperref}
\hypersetup{%
    breaklinks=true,%
	%pdfpagelabels=true,%
	plainpages=false,%
	pdfauthor={Author(s)},%
	pdftitle={\TITLE},%
	pdfsubject={\pdfSubject},%
	bookmarksnumbered=true,%
	colorlinks=false,%
	citecolor= black,%
	filecolor= black,%
	linkcolor=black,% you should probably change this to black before printing
	urlcolor=cyan,%
	pdfstartview=FitH%
}

\usepackage{float}
%\usepackage{indentfirst}
\usepackage{svg}
\usepackage{comment}

\usepackage{glossaries}
\makeglossaries
%%%%% Standard command %%%%%
% \gls{⟨label⟩}
 
%%%%% Capitalize first letter %%%%%
% \Gls{⟨label⟩}
 
%%%%% Pluralize term %%%%%
% \glspl{⟨label⟩}
 
%%%%% Capitalize and pluralize term %%%%%
% \Glspl{⟨label⟩}


\newglossaryentry{latex}
{
    name=latex,
    description={Is a markup language specially suited 
    for scientific documents}
}
\newglossaryentry{IC}{name= IC, description={Integrated Circuit}}
\newglossaryentry{RADAR}{name= RADAR, description={Radio Detection and Ranging}}
\newglossaryentry{HEMS}{name= HEMS, description={Helicopter Emergency Medical Service}}
\newglossaryentry{CSD}{name= CSD, description={Counter Surveliance Devices}}
\newglossaryentry{HPBW}{name= HPBW, description={Half Power Beam Width}}
\newglossaryentry{SHF}{name= SHF, description={Super High Frequency}}
\newglossaryentry{HLS}{name= HLS, description={High Level Specification}}
\newglossaryentry{DSP}{name= DSP, description={Digital Signal Processing}}
\newglossaryentry{RCS}{name= RCS, description={Radar Cross Section}}
\newglossaryentry{LNA}{name= LNA, description={Low Noise Amplifier}}
\newglossaryentry{CW}{name= CW, description={Continuous Wave}}
\newglossaryentry{AF}{name= AF, description={Array Factor}}
\newglossaryentry{NF}{name= NF, description={Near Field}}
\newglossaryentry{FF}{name= FF, description={Far Field}}
\newglossaryentry{AUT}{name= AUT, description={Antenna Under Test}}
\newglossaryentry{FFT}{name= FFT, description={Fast Fourier Transform}}
\newglossaryentry{IFFT}{name= IFFT, description={Inverse Fast Fourier Transform}}
\newglossaryentry{SNR}{name= SNR, description={Signal to Noise Ratio}}
\newglossaryentry{GNSS}{name = GNSS, description={Global Navigation Satellite System}}
\newglossaryentry{NF-FF}{name= NF-FF, description={Near Field to Far Field}}
\newglossaryentry{GPS}{name= GPS, description={Global Positioning System}}
\newglossaryentry{DGPS}{name= DGPS, description={Differential Global Positioning System}}
\newglossaryentry{PID}{name= PID, description={Proportional-Integral-Derivative}}
\newglossaryentry{RTK}{name= RTK, description={Real-Time Kinematic}}
\newglossaryentry{PPK}{name= PPK, description={Post-Processed Kinematic}}
\newglossaryentry{AAU}{name= AAU, description={Aalborg University}}
\newglossaryentry{RF}{name= RF, description={Radio frequency}}
\newglossaryentry{FNBW}{name=FNBW, description={First Null Beam Width}} 
\def\configurableCode{https://github.com/Mikkelnn/ESD5-Project/tree/main/spherical-NF-FF/configurable}
\def\testStandardGainHornDatasheet{https://github.com/Mikkelnn/ESD5-Project/blob/main/NF-FF-Data-2/16240-20\%20datasheet.pdf}
\def\nfTestDataFile{https://github.com/Mikkelnn/ESD5-Project/blob/main/NF-FF-Data-2/16240-20CBCFF_dir_30_010000.CSV}
\def\ffTestDataFile{https://github.com/Mikkelnn/ESD5-Project/blob/main/NF-FF-Data-2/Flann16240-20_CBC_010000.CSV}
\def\tesResultsGitHub{https://github.com/Mikkelnn/ESD5-Project/tree/main/spherical-NF-FF/testResults}
\def\systemValidationComparisonSmoothed{https://github.com/Mikkelnn/ESD5-Project/blob/main/Figures/systemDesign/validation/Comparison_plot_SNIFT_implementation_and_TICRA_data_(smoothed).svg}

\setlength\parindent{0pt}
